\section{Problema de Investigación}
Para plantear el problema que motiva este proyecto se describe primero la
situación actual, después su importancia, en seguida lo que aún desconocemos
de ella y, por último, la viabilidad de estudiarla e intervenirla.

\subsection{La Situación}
El ingreso de los estudiantes al Tecnológico de Estudios Superiores de Chalco
(TESCHA) se controla mediante la revisión visual de credenciales físicas por
parte del personal de vigilancia. Este método es lento en las horas de mayor
afluencia, no deja un registro digital de quién entra ni a qué hora y no
permite comprobar que quien presenta la credencial sea su titular. La
insatisfacción con los registros manuales se ha documentado en otros
contextos: en una escuela de educación básica de Perú, el 87.5\,\% de los 40
docentes y administrativos encuestados no se sentía cómodo con el registro
manual de asistencia \parencite{torres2019qr}. En el TESCHA, la magnitud de
estas limitaciones se medirá con los instrumentos de los Apéndices~A, B y~C.

El problema identificado es que el acceso de los estudiantes al TESCHA carece
de un mecanismo de validación automatizado. Como consecuencia, se forman
filas en los minutos previos al inicio de clases, una credencial prestada o
ajena puede superar la revisión visual y la institución no cuenta con un
registro confiable de ingresos que sirva para tomar decisiones.

Conviene distinguir este problema del pase de lista dentro del salón (la
práctica conocida como ``firmar por un compañero''). Ambos comparten la misma
raíz, una identificación deficiente, pero un registro en la entrada solo
demuestra que el estudiante ingresó al plantel, no que asistió a una clase
específica. Por ello, el control de asistencia por materia queda fuera del
alcance de este proyecto.

Esta situación se presenta en los accesos peatonales del TESCHA y se agudiza
en el turno vespertino, durante los 15 a 20 minutos previos al inicio de la
primera clase, cuando la mayor parte de los estudiantes llega en un intervalo
corto.

\subsection{La Importancia}
La importancia de este problema se aprecia en quiénes lo padecen:
\begin{enumerate}
	\item A los \textbf{estudiantes}, que pierden tiempo en las filas de
	      entrada, se exponen a llegar tarde a su primera clase y conviven en un
	      plantel donde el ingreso de personas ajenas no se controla con rigor.
	\item Al \textbf{personal de vigilancia}, que debe revisar cientos de
	      credenciales en pocos minutos sin ninguna herramienta que le ayude a
	      detectar una credencial ajena, vencida o apócrifa.
	\item A los \textbf{directivos y al Departamento de Control Escolar}, que no
	      disponen de datos de ingreso para detectar a estudiantes que dejan de acudir
	      al plantel ni para planear la apertura de accesos y el personal necesario en
	      cada horario.
\end{enumerate}

Y también en lo que provoca:
\begin{itemize}
	\item \textbf{Filas y tiempos muertos} al inicio de la jornada, que
	      provocan retardos en la primera clase. La teoría de colas, que ya se ha
	      propuesto para reducir los tiempos de espera de estudiantes en el área de
	      control escolar de otro plantel del Tecnológico Nacional de México
	      \parencite{talavera2025espera}, permite analizar este efecto a partir de la
	      tasa de llegadas y del tiempo de atención.
	\item \textbf{Riesgo de seguridad}, porque una credencial prestada, vencida
	      o ajena puede permitir el ingreso de personas que no pertenecen a la
	      comunidad del TESCHA.
	\item \textbf{Ausencia de datos}: sin un registro digital no es posible
	      saber cuántas personas ingresan ni en qué horarios, ni detectar a tiempo a
	      los estudiantes que dejan de acudir al plantel, una señal temprana de
	      posible deserción.
	\item \textbf{Carga operativa} para el personal de vigilancia, cuya
	      revisión depende únicamente de la atención de una persona en los momentos
	      de mayor afluencia.
\end{itemize}

\subsection{Lo que Desconocemos}
Para comprender mejor el problema se requiere información en cinco ejes.

\subsubsection{Situación Actual del Acceso en el TESCHA}
Es la información más importante, porque sin ella no es posible afirmar la
magnitud del problema. Se necesita conocer la matrícula por turno, el número
de accesos y de vigilantes, la tasa de llegada de estudiantes por minuto en el
horario pico, el tiempo que toma revisar a cada estudiante, la frecuencia con
que se olvidan o se prestan credenciales y la proporción de estudiantes que
cuenta con un celular propio. Estos datos se obtendrán con la encuesta del
Apéndice~A, la entrevista del Apéndice~B y la hoja de observación del
Apéndice~C.

\subsubsection{Arquitectura de Hardware y Tiempos de Respuesta}
\begin{itemize}
	\item \textbf{Latencia de lectura:} cuántos milisegundos toma al ESP32-CAM
	      capturar y decodificar el código, y cuántos más consume la consulta al
	      servidor por Wi-Fi. La suma debe ser menor que el tiempo de revisión manual
	      medido en el Apéndice~C para que el sistema tenga sentido. No se encontró
	      literatura arbitrada que reporte este tiempo en un ESP32-CAM, por lo que se
	      medirá de forma experimental con el prototipo; la referencia más cercana
	      disponible es un preprint que documenta el desempeño general de transmisión
	      de video del ESP32-CAM \parencite{nowroz2025esp32cam}, aunque no aborda la
	      decodificación de códigos QR.
	\item \textbf{Condiciones de lectura:} cómo afectan la iluminación, el
	      brillo de la pantalla del celular y la distancia al lector la tasa de
	      lecturas fallidas.
	\item \textbf{Durabilidad mecánica:} cuántos ciclos diarios deberá soportar
	      el mecanismo del torniquete según el flujo observado.
\end{itemize}

\subsubsection{Seguridad y Protección de Datos Personales}
\begin{itemize}
	\item \textbf{Validez del código:} qué duración del token equilibra la
	      seguridad frente a capturas de pantalla con la tolerancia a desfases de
	      reloj entre el celular y el servidor.
	\item \textbf{Comunicación segura:} uso de HTTPS entre el lector, el
	      servidor y el panel web para que los registros no puedan interceptarse ni
	      modificarse.
	\item \textbf{Marco legal:} como organismo público del Estado de México, el
	      TESCHA es sujeto obligado de la Ley de Protección de Datos Personales en
	      Posesión de Sujetos Obligados del Estado de México y Municipios
	      \parencite{edomex2017datos}, además de la Ley General de la materia
	      expedida en 2025 \parencite{lgpdppso2025}. Hay que identificar las
	      obligaciones que ambas imponen al sistema, como el aviso de privacidad y el
	      periodo de conservación de los registros de ingreso.
\end{itemize}

\subsubsection{Modelado de Datos e Indicadores}
\begin{itemize}
	\item \textbf{Estructura relacional:} diseño normalizado de las tablas de
	      estudiantes, accesos, dispositivos lectores y registros de ingreso en MySQL
	      o PostgreSQL.
	\item \textbf{Indicadores útiles:} qué reportes necesitan realmente
	      vigilancia y Control Escolar, por ejemplo, ingresos por franja horaria,
	      intentos rechazados y estudiantes sin ingresos durante varios días
	      consecutivos.
\end{itemize}

\subsubsection{Aceptación y Experiencia de Usuario}
\begin{itemize}
	\item \textbf{Disposición a usar el sistema:} si los estudiantes estarían
	      dispuestos a mostrar un QR en su celular y qué preocupaciones tienen sobre
	      el registro de su hora de entrada.
	\item \textbf{Ergonomía del lector:} altura y ángulo de la cámara para que
	      estudiantes de distintas estaturas lean su código sin detener el paso.
\end{itemize}

\subsection{La Viabilidad}
El uso de códigos QR para el control de acceso físico ya se ha propuesto en
la literatura \parencite{kao2011physical}, y su bajo costo frente a otras
tecnologías (Tabla~\ref{tab:comparativa}) hace viable construir y probar un
prototipo con los recursos del equipo. Frente a NFC, una comparación en
transacciones móviles encontró que el QR se procesa con mayor rapidez, aunque
no incorpora cifrado propio \parencite{harnaningrum2024comparison}; esa
carencia se compensa con la firma HMAC descrita en Delimitación. Además, el
estudiante usa su propio celular, por lo que no hace falta emitir tarjetas.
Frente al reconocimiento facial, el QR no requiere procesamiento pesado ni
trata datos biométricos, cuya aceptación entre estudiantes universitarios
depende en buena medida de la confianza y de la seguridad que perciben
\parencite{rukhiran2023user}, y cuyo tratamiento implicaría, además, las
obligaciones legales adicionales señaladas en ``Seguridad y Protección de
Datos Personales''.

\begin{table}[tbp]
	\caption{Comparación de Tecnologías de Validación de Acceso}
	\label{tab:comparativa}
	\begin{tabular}{lccc}
		\toprule
		Criterio                            & Código QR                & NFC     & Reconocimiento facial \\
		\midrule
		Costo de infraestructura            & Bajo                     & Medio   & Alto                  \\
		Medio que porta el estudiante       & Celular                  & Tarjeta & Ninguno               \\
		Complejidad de cómputo en el lector & Baja                     & Baja    & Alta                  \\
		Trata datos biométricos             & No                       & No      & Sí                    \\
		Riesgo de copia del identificador   & Medio\textsuperscript{a} & Bajo    & Bajo                  \\
		\bottomrule
	\end{tabular}
	\tablenote{Elaboración propia con base en
		\textcite{harnaningrum2024comparison} y \textcite{rukhiran2023user}. El costo
		y la complejidad son una valoración cualitativa del equipo.
		\textsuperscript{a}El riesgo corresponde a un QR estático; con un código
		dinámico firmado con HMAC se reduce considerablemente
		\parencite{subpratatsavee2015hmac}.}
\end{table}

Queda por conocer la viabilidad tecnológica desde el lado del estudiante, es
decir, si cuenta con un celular propio, con conexión a internet y con batería
suficiente al llegar al plantel; esto se indaga en la Sección~IV del
Apéndice~A.

La investigación es necesaria porque hoy no existen datos que midan el
problema en el TESCHA: se desconoce cuánto tiempo espera un estudiante para
entrar, con qué frecuencia se usan credenciales ajenas y si la comunidad
aceptaría un sistema basado en su celular. Levantar esa información y
contrastarla con el desempeño de un prototipo permitirá saber, con evidencia y
no con suposiciones, si esta solución de bajo costo ayuda a aliviar los
estragos del sistema actual y si justifica su adopción en el resto de los
accesos del plantel.
